% Generated draft: requires mathematical and source review.
\documentclass[11pt]{article}
\usepackage[margin=25mm]{geometry}
\usepackage{fontspec}
\setmainfont{Latin Modern Roman}
\usepackage{amsmath,amssymb,mathtools,unicode-math}
\setmathfont{Latin Modern Math}
\usepackage{enumitem,needspace,xurl,hyperref}
\setlength{\parindent}{0pt}
\setlength{\parskip}{5pt}
\setlength{\emergencystretch}{3em}
\newcommand{\N}{\mathbb{N}}
\newcommand{\Z}{\mathbb{Z}}
\newcommand{\Q}{\mathbb{Q}}
\newcommand{\R}{\mathbb{R}}
\newcommand{\C}{\mathbb{C}}
\newcommand{\F}{\mathbb{F}}
\newcommand{\A}{\mathbb{A}}
\newcommand{\PP}{\mathbb{P}}
\newcommand{\E}{\mathbb{E}}
\newcommand{\eps}{\varepsilon}
\newcommand{\dd}{\,\mathrm{d}}
\begin{document}
\section*{UnsolvedMath ID: 30004236}
\textbf{Source problem code:} catalogue entry
\subsection{Definitions and mathematical statement}
The source uses \(\operatorname{aff}(S)=\Z^d\) without specifying the meaning of \(\operatorname{aff}\). It cannot mean the usual real affine hull: for \(d\geq 1\), a real affine subspace cannot equal \(\Z^d\). The following formulation is therefore conditional on interpreting it as the integer affine span.

For a nonempty finite set \(S\subseteq\Z^d\), define its integer affine span by
\[
\operatorname{aff}_{\Z}(S)=\left\{\sum_{s\in S} n_s s : n_s\in\Z,\ \sum_{s\in S}n_s=1\right\}.
\]
For each positive integer \(d\), let \(\mathcal P_d\) be the collection of full-dimensional lattice polytopes \(P\subseteq\R^d\), meaning convex hulls of finite subsets of \(\Z^d\) with real affine dimension \(d\), such that
\[
\operatorname{aff}_{\Z}(P\cap\Z^d)=\Z^d.
\]
For \(P\in\mathcal P_d\), set
\[
\operatorname{SR}(P)=\min\left\{|B|:B\subseteq P\cap\Z^d,\ \operatorname{aff}_{\Z}(B)=\Z^d\right\},
\qquad
\operatorname{SR}(d)=\max_{P\in\mathcal P_d}\operatorname{SR}(P).
\]
The supplied record reports a finite bound on \(\operatorname{SR}(P)\) depending only on \(d\); this reported bound would ensure that the displayed maximum exists. The original statement supplies the values
\[
\operatorname{SR}(1)=2,\qquad \operatorname{SR}(2)=3,\qquad \operatorname{SR}(3)=5.
\]
Subject to confirmation of the integer affine span interpretation, determine \(\operatorname{SR}(4)\) and the asymptotic order of \(\operatorname{SR}(d)\) as \(d\to\infty\). The bound, the three stated values, and the open status are supplied source claims, not independently verified assertions in this draft.
\subsection{Short English statement}
Under the explicitly conditional interpretation of affine span as integer affine span, find the worst-case minimum number of lattice points needed to generate the entire ambient affine lattice from a spanning lattice polytope. Determine this number in dimension four and its asymptotic order as the dimension grows. The source's span notation requires confirmation.
\subsection{Sources}
\begin{itemize}
\item Local UnsolvedMath registry: permanent ID \(30004236\), title \emph{Spanning Rank of Lattice Polytopes}, collection \emph{Oberwolfach Reports Open Problems}, exact source code \texttt{catalogue entry}. The record supplies proposed year 2019 and names no individual proposer.
\item Supplied report citation: \emph{Geometric, Algebraic, and Topological Combinatorics} (2020), \emph{Oberwolfach Reports}, pp.~2395--2472, DOI \texttt{10.4171/owr/2019/39}. \url{https://doi.org/10.4171/owr/2019/39}. The background describes the report as documenting workshop open-problem discussions.
\item Supplied literature reference attributed to Averkov, Hofscheier, and Nill: \url{https://arxiv.org/abs/1911.03511}. The record provides no paper title or pinpoint reference.
\item Literature assessment dated exactly \texttt{2026-08-22}: the supplied curation labels the problem open, reports dimension-only boundedness and known values in dimensions one through three, and describes the dimension-four value and asymptotic order as unresolved. This is supplied dated curation, not an independently verified literature assessment. The record itself calls for independent verification of time-sensitive status and cited claims.
\item The registry labels its statement \texttt{corrected\_verified} and reports that the official Oberwolfach report was checked directly. These are supplied provenance claims; no source was independently visited or checked for this draft.
\end{itemize}
\end{document}
